\subsection{How the Church Learned to Keep Advent}\label{sec:history}

At Tours, by the tomb of St Martin, a bishop once laid down a fast before
Christmas. In the tenth book of his \textit{History of the Franks}, Gregory of
Tours tells of Perpetuus, the sixth bishop of his see, who built a larger
church for St Martin's body and laid down fasts and vigils for the whole round
of the year, a rule which, Gregory adds, \latin{hodieque apud nos tenetur
scriptum}, ``is kept among us in writing to this day''. Then he copies it out.
One line of it reads:

\begin{quote}
From the burial of our lord Martin to the Nativity of our Lord, three fasts each
week.\par
{\smallskip\noindent\small ---St Gregory of Tours, History of the Franks, Book X\par}
\end{quote}

Late in the sixth century a council of bishops at Mâcon, summoned by King
Guntrum, made the same rule a canon:

\begin{quote}
From S. Martin's Day until Christmas, a fast must be kept on Monday,
Wednesday, and Friday of each week\dots\ The Sacrifice must be offered after
the manner of the Lenten season.\par
{\smallskip\noindent\small ---First Council of Mâcon, canon 9\par}
\end{quote}

Dom Prosper Guéranger, reading the rule of Perpetuus in his \textit{Liturgical
Year}, counted forty-three days in it, ``consecrated to penance, as though it
were a second Lent, though less strict and severe than that which precedes
Easter.'' A council at Tours in 567 bound its monks to fast every day in
December until Christmas. Earlier still, a synod at Saragossa in Spain, about
the year 380, had required every one to come to church daily from 17 December
to the Epiphany. It asked for presence, not for a fast, and its first day is
the day on which the Roman Office now begins the great O antiphons.

\subsubsection*{The fast of the tenth month}

Rome came to the season by another road. Pope St Leo the Great preached a
series of sermons on a fast his Church kept in December, the tenth month, one of
four fasts set at the four seasons:

\begin{quote}
Accordingly we keep the spring fast in Lent, the summer fast at Whitsuntide, the
autumn fast in the seventh month, and the winter fast in this which is the tenth
month, knowing that there is nothing unconnected with the Divine commands, and
that all the elements serve the Word of God to our instruction, so that from the
very hinges on which the world turns, as if by four gospels we learn unceasingly
what to preach and what to do.\par
{\smallskip\noindent\small ---St Leo the Great, Sermon 19\par}
\end{quote}

Leo ends his sermons on this fast by naming its days: Wednesday and Friday, and a
vigil on Saturday with the Apostle Peter. They are the Ember Days of December,
which the Church still keeps in the third week of Advent, and the Saturday's
station is still St Peter's. Leo calls it simply the fast of the tenth month.

Around the fast grew the Sunday Masses of expectation, and the
books did not agree on how many there should be: the oldest Roman books have
five Masses for Advent or four Sundays, the Ambrosian rite of Milan has six, and
the old Spanish books never more than five. Amalarius of Metz, who wrote an
exposition of the offices and the Mass, found a fifth week before Christmas
marked in old Missals and lectionaries, and four weeks in the Missal he
ascribed, with a careful word, to St Gregory the Great: \latin{cujus auctorem
credimus esse beatum papam Gregorium}, ``whose author we believe to be the
blessed Pope Gregory''. Gregory's own sermon on the signs of the end, which the
Breviary reads on the First Sunday, was preached, its heading says, on the
Second Sunday of Advent.

\subsubsection*{What the season leaves unsung}

The medieval writers on the Church's offices knew Advent first by its silences,
and they asked what the silences meant. Amalarius had seen the \latin{Gloria in
excelsis}, Glory to God in the highest, left out of the Mass in the days of
Advent. Honorius Augustodunensis read the four weeks, with their silent
\latin{Gloria}, as the time of the Old Testament, which ran its course under
four orders of men who foreshadowed the Nativity: patriarchs, judges, kings and
prophets. To him each December was a passage through the centuries of waiting.

John Beleth, in the twelfth century, set down what else fell silent. In Advent,
as in Septuagesima, certain songs of gladness were left out: the \latin{Te Deum
laudamus}, We praise thee, O God; the \latin{Gloria in excelsis}; and the
dismissal at the end of Mass, \latin{Ite, missa est}. For the \latin{Te Deum} he
had two reasons. It was silent so that it might be taken up again with greater
ardour on the feast itself; or because \latin{nondum præsens est quem
exspectamus}, he whom we await is not yet here, and this is a hymn sung to one
who is present.

Yet the \latin{Alleluia} stayed, and Beleth explains why. In Septuagesima, he says,
the Church silences it because that whole season is a time of mourning; in
Advent she keeps it, because this time is partly a time of joy. William
Durandus, bishop of Mende, put the question sharply. If Advent is a fast of
mourning, why is the \latin{Alleluia} sung? If it is a fast of joy, why not the
\latin{Gloria} as well? His solution is that this fast is \latin{partim
exultationis, \& partim mœroris}, partly of exultation and partly of mourning.
The \latin{Gloria}, he says, is held back so that it may be sung the more
eagerly on the night of the Nativity.

These are the two roots of the season. From Gaul came weeks of penance, a
second and gentler Lent between St Martin and Christmas. From Rome came the
December fast and the Sunday Masses of expectation. The Advent of the Roman
books holds both: weeks of penance, with the \latin{Alleluia} still sounding
through them.
